\chapter{🔄 〜ている — Progreso y Estado}
\unidadheader{15}{〜ている}{Progreso y Estado}{Expresar acciones en curso y estados resultantes}

\begin{tcolorbox}[vocabbox, title={📌 Vocabulario Nuevo}]
\begin{tabularx}{\linewidth}{llX}
\toprule
\textbf{Japonés} & \textbf{Español} & \textbf{Tipo} \\
\midrule
\vocabitem{着る}{きる}{ponerse (ropa)}{v. gr.2}
\vocabitem{住む}{すむ}{vivir / residir en}{v. gr.1}
\vocabitem{知る}{しる}{saber / conocer}{v. gr.1}
\vocabitem{結婚する}{けっこんする}{casarse}{v. irr.}
\vocabitem{太る}{ふとる}{engordar}{v. gr.1}
\vocabitem{痩せる}{やせる}{adelgazar}{v. gr.2}
\vocabitem{疲れる}{つかれる}{cansarse}{v. gr.2}
\vocabitem{終わる}{おわる}{terminar (algo)}{v. gr.1}
\vocabitem{始まる}{はじまる}{empezar (algo intrans.)}{v. gr.1}
\vocabitem{慣れる}{なれる}{acostumbrarse}{v. gr.2}
\bottomrule
\end{tabularx}
\end{tcolorbox}

\subsection*{🗣️ Frases Clave}

\frase{いま\ruby{何}{なに}をしていますか？}{¿Qué estás haciendo ahora?}
\frase{テレビを\ruby{見}{み}ています。}{Estoy viendo la televisión.}
\frase{\ruby{東京}{とうきょう}に\ruby{住}{す}んでいます。}{Vivo en Tokio.}
\frase{\ruby{彼}{かれ}はもう\ruby{結婚}{けっこん}していますか？}{¿Él ya está casado?}

\subsection*{⚙️ Gramática}

\patron{Acción en progreso}
  {V-て + います (está haciendo ~)}
  {\ej{いま\ruby{電話}{でんわ}しています。}{Estoy hablando por teléfono ahora.}
  \ej{\ruby{雨}{あめ}が\ruby{降}{ふ}っています。}{Está lloviendo.}}

\patron{Estado resultante de una acción}
  {V-て + います (ya está en ese estado)}
  {Para verbos de cambio (結婚する、住む、知る、着る), ている = estado actual.
  \ej{\ruby{大阪}{おおさか}に\ruby{住}{す}んでいます。}{Vivo en Osaka. (estado)}
  \ej{\ruby{田中}{たなか}さんを\ruby{知}{し}っていますか？}
    {¿Conoce/sabe quién es el Sr. Tanaka?}
  \ej{はい、\ruby{知}{し}っています。}{Sí, lo conozco.}}

\patron{Negación: no está haciendo / no sabe}
  {V-て + いません}
  {\ej{まだ\ruby{食}{た}べていません。}{Aún no he comido.}
  \ej{\ruby{知}{し}りません。}{No lo sé. ← ojo: 知りません, no 知っていません}}

\begin{tcolorbox}[ejerciciobox, title={📝 Ejercicios Rápidos}]
\begin{enumerate}
  \item ¿Qué estás haciendo ahora? Responde con ています.
  \item Traduce: \ruby{彼女}{かのじょ}はもう\ruby{日本}{にほん}に\ruby{住}{す}んでいます。
  \item ¿Cuándo usas 知りません vs 知っていません?
\end{enumerate}
\vspace{0.4em}
\textbf{Respuestas:}
1) (personal)\quad
2) Ella ya vive en Japón.\quad
3) 知りません = respuesta estándar para "no sé"; 知っていません es técnicamente válido pero menos natural.
\end{tcolorbox}

\begin{tcolorbox}[kanjibox, title={🀄 Kanjis de la Unidad}]
\begin{tabularx}{\linewidth}{cllXX}
\toprule
\textbf{Kanji} & \textbf{On} & \textbf{Kun} & \textbf{Significado} & \textbf{Vocab} \\
\midrule
\kanjirow{着}{ちゃく}{き・つ}{llegar / ponerse ropa}{\ruby{着}{き}る / \ruby{到着}{とうちゃく}}
\kanjirow{住}{じゅう}{す}{vivir / residir}{\ruby{住}{す}む / \ruby{住所}{じゅうしょ}}
\kanjirow{知}{ち}{し}{saber / conocer}{\ruby{知}{し}る / \ruby{知識}{ちしき}}
\kanjirow{疲}{ひ}{つか}{cansarse}{\ruby{疲}{つか}れる / \ruby{疲労}{ひろう}}
\kanjirow{終}{しゅう}{お}{terminar}{\ruby{終}{お}わる / \ruby{終了}{しゅうりょう}}
\bottomrule
\end{tabularx}
\end{tcolorbox}
